\section{Related Work}
\label{sec:related-work}

\paragraph{Time-series forecasting and the foundation-model era.}
Per-series statistical models (ARIMA, exponential smoothing, state-space)
dominated forecasting until global deep models such as
DeepAR~\citep{salinas2020deepar} and N-BEATS~\citep{oreshkin2020nbeats}
showed that one network trained over many series can outperform
per-series fits. Transformers were beaten by simple linear baselines on
these benchmarks until PatchTST~\citep{nie2023patchtst} introduced patch
tokenization: contiguous values are grouped into a single token, the
attention cost drops by the patch length, and the model gains a useful
local prior. TiDE~\citep{das2023tide} pursued an MLP-only line in
parallel, with a residual block over patched inputs that is structurally
close to the encoder we use here.
TimesFM~\citep{das2024timesfm} is the canonical decoder-only patched TS
foundation model: pretrained on roughly 100B time points, its zero-shot
accuracy is on par with supervised specialists on Monash, Darts and ETT.
The follow-ups extend the recipe along different axes:
Chronos~\citep{ansari2024chronos} tokenizes values into a discrete
vocabulary and trains a T5 encoder-decoder over them;
MOIRAI~\citep{woo2024moirai} masks patches and reconstructs them with an
encoder; Time-MoE~\citep{shi2024timemoe} scales to a billion parameters
with sparse routing; Sundial~\citep{liu2025sundial} replaces categorical
heads with continuous flow matching;
Timer-S1~\citep{liu2026timers1} currently leads GIFT-Eval. All of them
forecast in value space, whether the head is regressive, categorical or
generative.

\paragraph{Contrastive representation learning for time series.}
A separate line learns time-series representations by discrimination
rather than reconstruction. Contrastive Predictive
Coding~\citep{oord2018cpc} encodes a series with a strided CNN,
summarises the past with a GRU, and trains the encoder to assign higher
score to the true future latent than to a batch of negatives. The loss,
InfoNCE~\citep{oord2018cpc}, is a tractable lower bound on the mutual
information between the context and the future, so the encoder is
forced to keep whatever in the past predicts the future. The
time-series successors keep the contrastive backbone but rework what
counts as a positive pair: subseries consistency in
T-Loss~\citep{franceschi2019tloss}, temporal neighbourhoods in
TNC~\citep{tonekaboni2021tnc}, contextual consistency under masking and
random crops with hierarchical contrast in
TS2Vec~\citep{yue2022ts2vec}. These methods produce strong
representations for classification and anomaly detection on UCR and UEA,
but no current TS foundation model uses contrastive learning as the
primary pretraining objective.

\paragraph{Latent-space prediction without contrast.}
A third line, originating in vision, predicts in latent space rather
than in input space. The Joint Embedding Predictive
Architecture~\citep{assran2023ijepa} encodes a context block with one
network, encodes the target block with a second network whose weights
are an exponential moving average of the first, and trains a predictor
to map the context latent to the target latent. The loss is L2 in
latent space. Collapse is avoided by the combination of the asymmetric predictor and
the stop-gradient on the EMA target rather than by contrast: there are
no negatives. Two recent papers
carry the idea to time series. LaT-PFN~\citep{verdenius2024latpfn}
combines JEPA with a prior-data fitted network for in-context
forecasting, and TS-JEPA~\citep{ennadir2025tsjepa} adapts the
patch-token JEPA recipe to univariate series, with competitive results
on classification and long-term forecasting. Masked-reconstruction TS
models such as MOMENT~\citep{goswami2024moment} and
MOIRAI~\citep{woo2024moirai} sit closer to JEPA than tokenized or
autoregressive lines do, since they too predict masked content from
visible context, but their target remains the raw values rather than
their representation.

\paragraph{Positioning.}
Decoder-only transformers trained to predict the next token have become
the dominant architecture in language. As Sutskever has argued, this
objective is far less narrow than it looks: predicting the next token
correctly requires the model to encode grammar, style, world knowledge
and arithmetic, since each is a sub-task of the same prediction
problem. We want to recover this property for time series, with a single
decoder-only model that predicts the next patch and learns whatever
structure is useful along the way, but without committing to
value-space prediction. The observed series is a sample from a
distribution generated by a stochastic process, so a value-space loss
penalises the model for failing to predict intrinsic noise that no
model can forecast. Predicting in latent space lets the target be a
representation of the distribution of plausible futures, with the
observed value treated as one of its samples, and the model is not
forced to fit details that carry no semantic content.

Contrastive forecasting is our attempt at this combination. We forecast
the next patch in latent space, as in JEPA, but enforce non-collapse
with explicit InfoNCE~\citep{oord2018cpc} negatives, as in CPC and
TS2Vec. JEPA is a siamese architecture and is therefore at risk of
collapse by construction: the predictor can drive the L2 loss to zero
by mapping every context to the same point. With an InfoNCE loss the
constant function is instead one of the least stable solutions, since
every sample is pushed away from every other sample in latent space,
and the collapsed embedding is actively repelled. Compared to CPC and
its successors, the target is forecasting rather than representation
transfer for classification. To our knowledge, no published work has
measured latent-space forecasting with explicit negatives at TS-FM
scale.
